\documentclass[a4paper, 11pt]{article}
\usepackage{graphicx} % Required for inserting images

% pseudocode 
\usepackage{algorithm}
\usepackage{algpseudocode}
\usepackage{underscore}

\usepackage{enumitem}

\usepackage{hyperref}

\usepackage{amssymb, amsfonts, amsmath, amsthm}

\makeatletter
\newcommand\xlongrightarrow[2][]{%
  \ext@arrow 9999{\longleftrightarrowfill@}{#1}{#2}}
\newcommand\longrightarrowfill@{%
  \arrowfill@\leftarrow\relbar\rightarrow}
\makeatother

\theoremstyle{plain} 
\newtheorem{theorem}[equation]{Theorem} 
\newtheorem{corollary}[equation]{Corollary} 
\newtheorem{lemma}[equation]{Lemma} 
\newtheorem{proposition}[equation]{Proposition}

\theoremstyle{definition} 
\newtheorem{definition}[equation]{Definition} 

\theoremstyle{remark} 
\newtheorem{remark}[equation]{Remark} 
\newtheorem{ex}[equation]{Example} 
\newtheorem{notation}[equation]{Notation} 
\newtheorem{terminology}[equation]{Terminology}

\title{Reduced Order Waveform Relaxation}
\author{Anar Abdullayev\\[2ex]\textit{Faculty of Mathematics and Computer Science, University of Münster}}
\date{\today}

\usepackage{geometry}
\geometry{
a4paper, 
left = 1.5cm,
right = 1.5cm,
top = 2.54cm,
bottom = 2.54cm
}

\begin{document}

\maketitle

\begin{abstract}
    aim is to show things work as predicted...
\end{abstract}

\section{Introduction}
We consider the heat equation on a bounded domain $\Omega\subset \mathbb{R}^{n}$ with a smooth boundary $\partial\Omega$
\begin{align*}
    \dfrac{\partial u}{\partial t} &= \Delta u + f(x, t), \qquad x\in \Omega, \,\,\,\,\quad 0<t<T,\\
    u(x, t) &= g(x, t), \hspace{52pt} x\in \partial\Omega, \quad 0<t<T,\\
    u(x, 0) &= u_{0}(x), \hspace{55pt} x\in \Omega,
\end{align*}
where the initial condition $u_{0}(x)$ and the boundary condition $g(x, t)$ are bounded piecewise continuous and $f(x, t)$ is continuous. On each subdomain $\Omega_{j}$, we solve at each step $k+1$ of the waveform relaxation iteration the subproblem
\begin{align*}
    \dfrac{\partial u_{j}^{k+1}}{\partial t} &= \Delta u_{j}^{k+1} + f(x, t),\\
    u_{j}^{k+1}(x, t) &= u_{l}^{k}(x, t), \qquad x\in \Gamma_{jl}, \quad 0<t<T,\\
    u_{j}^{k+1}(x, t) &= g(x, t), \qquad x\in \Gamma_{i0}, \quad 0<t<T,\\
    u_{j}^{k+1}(x, 0) &= u_{0}(x), \qquad x\in \Omega_{j},
\end{align*}
for $j=1, \ldots, N$. Let us define the transmission operator $\mathcal{T}_{jl}: V_l \to V_{l}|_{\Gamma_{jl}}$ by
\begin{equation*}
    \mathcal{T}_{jl}u_{l}^{k} := \left.u_{l}^{k}\right|_{\Gamma_{jl}},
\end{equation*}
where $u_{l}^k \in V_{l}$ is the solution on subdomain $l$ at iteration $k$, $V_{l}$ denotes the function space associated with subdomain $l$, and $\Gamma_{jl}$ is the interface between subdomains $j$ and $l$. This operator extracts the values of the solution $u_{l}^k$ on the interface $\Gamma_{jl}$ between subdomains $j$ and $l$, which are then used as boundary data for subdomain $j$ in the overlapping waveform relaxation iteration.

Let $V_{j, r}\subset V_{j}$ be the reduced space, constructed from snapshots of $u_{j}$ or a global solution approximation. On each subdomain $\Omega_{j}$, we solve at each step $k+1$ of the reduced-order waveform relaxation iteration the subproblem
\begin{align*}
    \left(\partial_{t} u_{j, r}^{k+1}, v_{r}\right)_{\Omega_{j}} &=  \left(\Delta u_{j, r}^{k+1}, v_{r}\right)_{\Omega_{j}} + \left(f, v_{r}\right)_{\Omega_{j}}, \hspace{-50pt}&& \forall v_{r}\in V_{j, r},\\
    u_{j, r}^{k+1}(x, t) &= u_{l, r}^{k}(x, t), && x\in \Gamma_{jl}, \quad 0<t<T,\\
    u_{j, r}^{k+1}(x, t) &= P_{j, r}g(x, t), && x\in \Gamma_{j0}, \quad 0<t<T,\\
    u_{j, r}^{k+1}(x, 0) &= P_{j, r}u_{0}(x), && x\in \Omega_{j},
\end{align*}
where the projection operator $P_{j, r}: V_{j} \to V_{j, r}$ is defined as
\begin{equation*}
    P_{j,r} u := \arg\min_{v_r \in V_{j,r}} \lVert u - v_r \rVert_{V_j},
\end{equation*}
i.e., for any function $u \in V_j$, $P_{j,r} u$ is the best approximation of $u$ in the reduced space $V_{j,r}$ with respect to the norm $\|\cdot\|_{V_j}$. We note that the transmission operator defined above can also be applied to reduced-order solutions $u_{l,r}^k \in V_{l,r} \subset V_l$, so that the reduced interface condition in the subproblem can be written compactly as
\begin{equation*}
    \left.u_{j,r}^{k+1}\right|_{\Gamma_{jl}} = \mathcal{T}_{jl} u_{l,r}^{k}.
\end{equation*}
Now, let us define the error
\begin{equation*}
    e_{j, r}^{k+1} : = u_{j}^{k+1} - u_{j, r}^{k+1}\in V_{j}.
\end{equation*}
We can split this error as
\begin{equation*}
    e_{j, r}^{k+1} = \left(u_{j}^{k+1} - P_{j,r}u_{j}^{k+1}\right) + \left(P_{j,r}u_{j}^{k+1} - u_{j, r}^{k+1}\right).
\end{equation*}
Let us introduce the following
\begin{equation*}
    \eta_{j}^{k+1} = u_{j}^{k+1} - P_{j,r}u_{j}^{k+1}, \qquad \mu_{j}^{k+1} = P_{j,r}u_{j}^{k+1} - u_{j, r}^{k+1},
\end{equation*}
then the error becomes
\begin{equation*}
    e_{j, r}^{k+1} = \eta_{j}^{k+1} + \mu_{j}^{k+1}.
\end{equation*}
In the original error $e_{j, r}^{k+1}$, we are comparing functions in different spaces. That's why we split in such a way so that we can define error in a compatible way. 

Let us concentrate on the error $\mu_{j}^{k+1}$. It is clear that $P_{j,r}u_{j}^{k+1}\in V_{j, r}$. Then the projected problem reads
\begin{equation*}
    \left(\partial_{t} P_{j,r}u_{j}^{k+1}, v_{r}\right)_{\Omega_{j}} =  \left(\Delta P_{j,r}u_{j}^{k+1}, v_{r}\right)_{\Omega_{j}} + \left(f, v_{r}\right)_{\Omega_{j}} - \left(\Delta \left(P_{j,r}u_{j}^{k+1} - u_{j}^{k+1}\right), v_{r}\right)_{\Omega_{j}}, \quad \forall v_{r}\in V_{j, r},
\end{equation*}
where $\Delta \left(P_{j,r}u_{j}^{k+1} - u_{j}^{k+1}\right)$ represents the residual caused by the projection. Subtracting the reduced-order waveform
relaxation equation from this equation, we obtain
\begin{equation*}
    \left(\partial_{t} P_{j,r}u_{j}^{k+1} - \partial_{t} u_{j, r}^{k+1}, v_{r}\right)_{\Omega_{j}} =  \left(\Delta \left(P_{j,r}u_{j}^{k+1} - u_{j, r}^{k+1}\right), v_{r}\right)_{\Omega_{j}} - \left(\Delta \left(P_{j,r}u_{j}^{k+1} - u_{j}^{k+1}\right), v_{r}\right)_{\Omega_{j}}, \forall v_{r}\in V_{j, r}.
\end{equation*}
Considering the defined error definitions above, this reduces to
\begin{equation*}
    \left(\partial_{t}\mu_{j}^{k+1}, v_{r}\right)_{\Omega_{j}} =  \left(\Delta \mu_{j}^{k+1}, v_{r}\right)_{\Omega_{j}} + \left(\Delta \eta_{j}^{k+1}, v_{r}\right)_{\Omega_{j}}, \qquad \forall v_{r}\in V_{j, r}.
\end{equation*}
Now, let us consider transmission conditions for $\mu_{j}^{k+1}$. Since 
\begin{equation*}
    \left.u_{j}^{k+1}\right|_{\Gamma_{jl}} = \mathcal{T}_{jl}u_{l}^{k}, \qquad \left.u_{j, r}^{k+1}\right|_{\Gamma_{jl}} = \mathcal{T}_{jl}u_{l, r}^{k},
\end{equation*}
on the interface boundary $\Gamma_{jl}$, we have
\begin{align*}
    \left.\mu_{j}^{k+1}\right|_{\Gamma_{jl}} &= \left.\left(P_{j,r}u_{j}^{k+1} - u_{j, r}^{k+1}\right)\right|_{\Gamma_{jl}} = P_{j,r}\mathcal{T}_{jl}u_{l}^{k} - \mathcal{T}_{jl}u_{l, r}^{k} = \mathcal{T}_{jl}\mu_{l}^{k} + P_{j, r}\mathcal{T}_{jl}u_{l}^{k} - \mathcal{T}_{jl}P_{l, r}u_{l}^{k} = \\
    &= \mathcal{T}_{jl}\mu_{l}^{k} + T_{jl}\eta_{l}^{k} + P_{j, r}\mathcal{T}_{jl}u_{l}^{k} - \mathcal{T}_{jl}u_{l}^{k} = \mathcal{T}_{jl}\mu_{l}^{k} + \mathcal{T}_{jl}\eta_{l}^{k} + \left(P_{j, r} -I \right)\mathcal{T}_{jl}u_{l}^{k} =\\
    &= \mathcal{T}_{jl}\left(\mu_{l}^{k} + \eta_{l}^{k}\right) + \left(P_{j, r} -I \right)\mathcal{T}_{jl}u_{l}^{k} = \mathcal{T}_{jl}e_{l, r}^{k} + \left(P_{j, r} -I \right)\mathcal{T}_{jl}u_{l}^{k}.
\end{align*}
\begin{remark}
    The equality
    \begin{equation*}
        \left.\mu_{j}^{k+1}\right|_{\Gamma_{jl}} = \mathcal{T}_{jl}e_{l, r}^{k} + \left(P_{j, r} -I \right)\mathcal{T}_{jl}u_{l}^{k}
    \end{equation*}
    expresses how the interface error in the reduced problem is formed. One part of the interface error is simply the neighboring subdomain’s total error transported through the transmission operator, which is exactly the same mechanism by which errors propagate in the full waveform relaxation method and has nothing to do with model reduction. The other part is a purely reduction-induced defect: even if the neighboring solution were exact, projecting the transmitted interface data into the reduced space generally introduces a mismatch because the reduced space cannot represent all possible traces. Our reduced space $V_{j, r}$ is built from a finite number of snapshots of solutions inside the subdomain. These snapshots capture the dominant interior behavior of the solution, but they do not span all functions that can appear on the interface $\Gamma_{jl}$. In particular, the interface data coming from a neighboring subdomain through the transmission operator may contain spatial or temporal features that were never present in the snapshot set. When such interface data is imposed on the reduced subproblem, it must first be projected onto $V_{j, r}$, and during this projection, any component that lies outside the span of the reduced basis is inevitably lost. \textit{This loss shows up as a mismatch between the exact transmitted interface value and its reduced representation.} Therefore, even if the neighboring solution were exact, the reduced subdomain cannot reproduce its trace perfectly unless the reduced space is rich enough to contain that trace. This is precisely what is meant by saying that the reduced space cannot represent all possible traces, and it is why the additional projection term appears in the interface condition for the reduced error. Thus, the interface value of the reduced error is the sum of the propagated waveform relaxation error and a fixed consistency error caused by the projection. This separation is fundamental, because it shows that the iterative behavior of the reduced method mirrors that of the full method, while the influence of model reduction enters only through an additional, clearly identifiable projection error term.
\end{remark}
\noindent On the interface boundary $\Gamma_{j0}$, we have
\begin{equation*}
    \left.\mu_{j}^{k+1}\right|_{\Gamma_{j0}} = \left.\left(P_{j,r}u_{j}^{k+1} - u_{j, r}^{k+1}\right)\right|_{\Gamma_{j0}} = P_{j, r}g(x, t) - P_{j, r}g(x, t) = 0.
\end{equation*}
The initial condition becomes
\begin{equation*}
    \mu_{j}^{k+1}(x, 0) = 0.
\end{equation*}
Therefore, the problem for $\mu_{j}^{k+1}\in V_{j, r}$ is found as
\begin{align*}
    \left(\partial_{t}\mu_{j}^{k+1}, v_{r}\right)_{\Omega_{j}} &=  \left(\Delta \mu_{j}^{k+1}, v_{r}\right)_{\Omega_{j}} + \left(\Delta \eta_{j}^{k+1}, v_{r}\right)_{\Omega_{j}}, \hspace{-50pt}&& \forall v_{r}\in V_{j, r},\\
    \mu_{j}^{k+1}(x, t) &= \mathcal{T}_{jl}e_{l, r}^{k} + \left(P_{j, r} -I \right)\mathcal{T}_{jl}u_{l}^{k}, && x\in \Gamma_{jl}, \quad 0<t<T,\\
    \mu_{j}^{k+1}(x, t) &= 0, && x\in \Gamma_{j0}, \quad 0<t<T,\\
    \mu_{j}^{k+1}(x, 0) &= 0, && x\in \Omega_{j}.
\end{align*}
\begin{remark}
    As the reduced space $V_{j,r}$ is enriched, the projection operator $P_{j,r}$ converges to the identity on $V_j$. As a result, the projection error $\eta_j^{k+1} = u_j^{k+1} - P_{j,r}u_j^{k+1}$ diminishes, and the additional interface discrepancy $(P_{j,r} - I)\mathcal{T}_{jl}u_l^k$ becomes negligible. In this limit, the reduced error component $\mu_j^{k+1}$ increasingly coincides with the total error $e_{j,r}^{k+1}$, which is already evident from the decomposition $e_{j,r}^{k+1} = \eta_j^{k+1} + \mu_j^{k+1}$. Consequently, for sufficiently rich reduced spaces, the reduced-order waveform relaxation method reproduces the error propagation behavior of the full method, with the influence of model reduction becoming asymptotically negligible.
\end{remark}

\noindent Let us now choose
\begin{equation*}
    v_{r} = \mu_{j}^{k+1}\in V_{j, r}.
\end{equation*}
Then, the left-hand side becomes
\begin{equation*}
    \left(\partial_{t}\mu_{j}^{k+1}, \mu_{j}^{k+1}\right)_{\Omega_{j}} = \dfrac{1}{2} \dfrac{d}{dt} \left(\mu_{j}^{k+1}, \mu_{j}^{k+1}\right)_{\Omega_{j}} = \dfrac{1}{2}\dfrac{d}{dt}\left\lVert \mu_{j}^{k+1}\right\rVert^{2}_{L^{2}(\Omega_{j})}.
\end{equation*}
For the first term on the left-hand side, using Green’s first identity (integration by parts), we obtain
\begin{align*}
    \left(\Delta \mu_{j}^{k+1}, \mu_{j}^{k+1}\right)_{\Omega_{j}} &= \int_{\Omega_{j}}\Delta \mu_{j}^{k+1} \mu_{j}^{k+1} = -\int_{\Omega_{j}}\nabla \mu_{j}^{k+1} \cdot \nabla \mu_{j}^{k+1} + \int_{\partial\Omega_{j}}\mu_{j}^{k+1}\nabla\mu_{j}^{k+1}\cdot n=\\
    &= -\int_{\Omega_{j}}\nabla \mu_{j}^{k+1} \cdot \nabla \mu_{j}^{k+1} + \sum_{l\in \mathcal{N}_{j}}\int_{\Gamma_{jl}}\mu_{j}^{k+1}\nabla\mu_{j}^{k+1}\cdot n=\\
    &= -\lVert \nabla \mu_{j}^{k+1}\rVert^{2}_{L^{2}(\Omega_{j})} + \sum_{l\in \mathcal{N}_{j}}\int_{\Gamma_{jl}}\mu_{j}^{k+1}\nabla\mu_{j}^{k+1}\cdot n,
\end{align*}
where the neighbor index set is defined as
\begin{equation*}
    \mathcal{N}_{i} = \{j\in \{1, \ldots, N\}: \left(\partial\Omega_{i}\cap\Omega_{j}\right)\setminus \partial\Omega \neq \emptyset\}.
\end{equation*}
Since $\left.\mu_{j}^{k+1}\right|_{\Gamma_{jl}} = \mathcal{T}_{jl}e_{l, r}^{k} + \left(P_{j, r} -I \right)\mathcal{T}_{jl}u_{l}^{k}$, the above becomes
\begin{equation*}
    \left(\Delta \mu_{j}^{k+1}, \mu_{j}^{k+1}\right)_{\Omega_{j}} = -\lVert \nabla \mu_{j}^{k+1}\rVert^{2}_{L^{2}(\Omega_{j})} + \sum_{l\in \mathcal{N}_{j}}\int_{\Gamma_{jl}}\left(\mathcal{T}_{jl}e_{l, r}^{k} + \left(P_{j, r} -I \right)\mathcal{T}_{jl}u_{l}^{k}\right)\nabla\mu_{j}^{k+1}\cdot n.
\end{equation*}
Thus, the general equation becomes
\begin{equation*}
    \dfrac{1}{2}\dfrac{d}{dt}\left\lVert \mu_{j}^{k+1}\right\rVert^{2}_{L^{2}(\Omega_{j})} + \lVert \nabla \mu_{j}^{k+1}\rVert^{2}_{L^{2}(\Omega_{j})} = \sum_{l\in \mathcal{N}_{j}}\int_{\Gamma_{jl}}\left(\mathcal{T}_{jl}e_{l, r}^{k} + \left(P_{j, r} -I \right)\mathcal{T}_{jl}u_{l}^{k}\right)\nabla\mu_{j}^{k+1}\cdot n + \left(\Delta \eta_{j}^{k+1}, v_{r}\right)_{\Omega_{j}}.
\end{equation*}
Since $\mu_{j}^{k+1}\in V_{j, r}\subset V_{j}\subset H^{1}(\Omega_{j})$, then its trace $\left.\mu_{j}^{k+1}\right|_{\partial\Omega_{j}}\in H^{1/2}(\partial\Omega_{j})$, and the (weak) normal derivative $\nabla\mu_{j}^{k+1}\cdot n$ is well-defined as an element of the dual space $H^{-1/2}(\partial\Omega_{j})$, in particular on the interface $\Gamma_{jl}\subset \partial \Omega_{j}$. In other words, $g^{k+1}_{j} = \mathcal{T}_{jl}e_{l, r}^{k} + \left(P_{j, r} -I \right)\mathcal{T}_{jl}u_{l}^{k} \in H^{1/2}(\Gamma_{jl})$. Using the Cauchy-Schwarz inequality, we get
\begin{equation*}
    \left|\int_{\Gamma_{jl}}g^{k+1}_{j}\nabla\mu_{j}^{k+1}\cdot n\right| \leq \lVert \mathcal{T}_{jl}e_{l, r}^{k} + \left(P_{j, r} -I \right)\mathcal{T}_{jl}u_{l}^{k} \rVert_{H^{1/2}(\Gamma_{jl})}\lVert \nabla\mu_{j}^{k+1}\cdot n \rVert_{H^{-1/2}(\Gamma_{jl})}.
\end{equation*}
Before proceeding, note that $g^{k+1}_{j}$ comes from subdomain $\Omega_{l}$ and is only defined on the interface. Let us consider the first term on the right-hand side:
\begin{equation*}
    \lVert \mathcal{T}_{jl}e_{l, r}^{k} + \left(P_{j, r} -I \right)\mathcal{T}_{jl}u_{l}^{k}\rVert_{H^{1/2}(\Gamma_{jl})} \leq \lVert \mathcal{T}_{jl}e_{l, r}^{k}\rVert_{H^{1/2}(\Gamma_{jl})} +  \lVert\left(P_{j, r} -I \right)\mathcal{T}_{jl}u_{l}^{k} \rVert_{H^{1/2}(\Gamma_{jl})}.
\end{equation*}
Using the note stated above and the interior trace\footnote{In a 2-dimensional domain, interior trace means curve inside the domain} inequality, we have
\begin{align*}
     \lVert \mathcal{T}_{jl}e_{l, r}^{k} \rVert_{H^{1/2}(\Gamma_{jl})} \leq C_{1, l} \lVert e_{l, r}^{k}\rVert_{H^{1}(\Omega_{l})},
\end{align*}
and the standard trace inequality
\begin{equation*}
    \lVert \nabla\mu_{j}^{k+1} \cdot n\rVert_{H^{-1/2}(\Gamma_{jl})} \leq C_{2, l}\lVert \nabla\mu_{j}^{k+1} \rVert_{L^{2}(\Omega_{j})}.
\end{equation*}
The interface defect is exactly the distance between the transmitted exact interface data and the reduced trace space
\begin{equation*}
    \lVert \left(P_{j, r} -I \right)\mathcal{T}_{jl}u_{l}^{k} \rVert_{H^{1/2}(\Gamma_{jl})} = \inf_{v_{r}\in V_{j, r}}\lVert \mathcal{T}_{jl}u_{l}^{k} -  v_{r} \rVert_{H^{1/2}(\Gamma_{jl})}.
\end{equation*}
Then, we obtain
\begin{equation*}
    \left|\int_{\Gamma_{jl}}g_{j}^{k+1}\nabla\mu_{j}^{k+1}\cdot n\right| \leq C_{2, l}\lVert \nabla\mu_{j}^{k+1} \rVert_{L^{2}(\Omega_{j})}\left(C_{1, l} \lVert e_{l, r}^{k}\rVert_{H^{1}(\Omega_{l})} + \inf_{v_{r}\in V_{j, r}}\lVert \mathcal{T}_{jl}u_{l}^{k} -  v_{r} \rVert_{H^{1/2}(\Gamma_{jl})}\right).
\end{equation*}
Let us focus on $\lVert \nabla\mu_{j}^{k+1} \rVert_{L^{2}(\Omega_{j})} \lVert e_{l, r}^{k}\rVert_{H^{1}(\Omega_{l})}$. Now, let us use Young's inequality to get
\begin{align*}
    \lVert e_{l, r}^{k} \rVert_{H^{1}(\Omega_{l})}\lVert \nabla\mu_{j}^{k+1} \rVert_{L^{2}(\Omega_{j})} &\leq \dfrac{\varepsilon}{2}\lVert e_{l, r}^{k} \rVert^{2}_{H^{1}(\Omega_{l})} + \dfrac{1}{2\varepsilon}\lVert \nabla\mu_{j}^{k+1} \rVert^{2}_{L^{2}(\Omega_{j})}, \quad \varepsilon > 0,\\
    \inf_{v_{r}\in V_{j, r}}\lVert \mathcal{T}_{jl}u_{l}^{k} -  v_{r} \rVert_{H^{1/2}(\Gamma_{jl})}\lVert \nabla\mu_{j}^{k+1} \rVert_{L^{2}(\Omega_{j})} &\leq \dfrac{\Tilde{\varepsilon}}{2}\inf_{v_{r}\in V_{j, r}}\lVert \mathcal{T}_{jl}u_{l}^{k} -  v_{r} \rVert^{2}_{H^{1/2}(\Gamma_{jl})} + \dfrac{1}{2\Tilde{\varepsilon}}\lVert \nabla\mu_{j}^{k+1} \rVert^{2}_{L^{2}(\Omega_{j})}, \quad \Tilde{\varepsilon} > 0.
\end{align*}
The inequality becomes
\begin{align*}
    \left|\int_{\Gamma_{jl}}g_{j}^{k+1}\nabla\mu_{j}^{k+1}\cdot n\right| &\leq C_{1, l}C_{2, l}\dfrac{\varepsilon}{2}\lVert e_{l, r}^{k} \rVert^{2}_{H^{1}(\Omega_{l})} + \left(C_{1, l}C_{2, l}\dfrac{1}{2\varepsilon} + C_{2, l}\dfrac{1}{2\Tilde{\varepsilon}}\right)\lVert \nabla\mu_{j}^{k+1} \rVert^{2}_{L^{2}(\Omega_{j})} + \\
    &\quad + C_{2, l}\dfrac{\Tilde{\varepsilon}}{2}\inf_{v_{r}\in V_{j, r}}\lVert \mathcal{T}_{jl}u_{l}^{k} -  v_{r} \rVert^{2}_{H^{1/2}(\Gamma_{jl})}.
\end{align*}
For the last term on the right-hand side, if we use the Cauchy-Schwarz inequality and Young's inequality, we get 
\begin{equation*}
    \left|\left(\Delta \eta_{j}^{k+1}, \mu_{j}^{k+1}\right)_{\Omega_{j}}\right| \leq \lVert \Delta\eta_{j}^{k+1} \rVert_{L^{2}(\Omega_{j})}\lVert \mu_{j}^{k+1} \rVert_{L^{2}(\Omega_{j})} \leq \dfrac{\hat{\varepsilon}}{2}\lVert \Delta\eta_{j}^{k+1} \rVert^{2}_{L^{2}(\Omega_{j})} + \dfrac{1}{2\hat{\varepsilon}}\lVert \mu_{j}^{k+1} \rVert^{2}_{L^{2}(\Omega_{j})}, \quad \hat{\varepsilon} > 0.
\end{equation*}
Considering all the inequalities derived for the right-hand side terms of the equation, we have
\begin{align*}
    \dfrac{1}{2}\dfrac{d}{dt}\left\lVert \mu_{j}^{k+1}\right\rVert^{2}_{L^{2}(\Omega_{j})} + \lVert \nabla \mu_{j}^{k+1}\rVert^{2}_{L^{2}(\Omega_{j})} &\leq \dfrac{\varepsilon}{2}\sum_{l\in \mathcal{N}_{j}}C_{1, l}C_{2, l}\lVert e_{l, r}^{k} \rVert^{2}_{H^{1}(\Omega_{l})} + \dfrac{\Tilde{\varepsilon}}{2}\sum_{l\in \mathcal{N}_{j}}C_{2, l}\inf_{v_{r}\in V_{j, r}}\lVert \mathcal{T}_{jl}u_{l}^{k} -  v_{r} \rVert_{H^{1/2}(\Gamma_{jl})}+\\
    &\quad + \sum_{l\in \mathcal{N}_{j}}\left(C_{1, l}C_{2, l}\dfrac{1}{2\varepsilon} + C_{2, l}\dfrac{1}{2\Tilde{\varepsilon}}\right)\lVert \nabla\mu_{j}^{k+1} \rVert^{2}_{L^{2}(\Omega_{j})}+\\
    &\quad + \dfrac{\hat{\varepsilon}}{2}\lVert \Delta\eta_{j}^{k+1} \rVert^{2}_{L^{2}(\Omega_{j})} + \dfrac{1}{2\hat{\varepsilon}}\lVert \mu_{j}^{k+1} \rVert^{2}_{L^{2}(\Omega_{j})},
\end{align*}
or equivalently,
\begin{align*}
    \dfrac{1}{2}\dfrac{d}{dt}\left\lVert \mu_{j}^{k+1}\right\rVert^{2}_{L^{2}(\Omega_{j})} &\leq \dfrac{\varepsilon}{2}\sum_{l\in \mathcal{N}_{j}}C_{1, l}C_{2, l}\lVert e_{l, r}^{k} \rVert^{2}_{H^{1}(\Omega_{l})} + \dfrac{\Tilde{\varepsilon}}{2}\sum_{l\in \mathcal{N}_{j}}C_{2, l}\inf_{v_{r}\in V_{j, r}}\lVert \mathcal{T}_{jl}u_{l}^{k} -  v_{r} \rVert_{H^{1/2}(\Gamma_{jl})}+\\
    &\quad + \left(\sum_{l\in \mathcal{N}_{j}}\left(C_{1, l}C_{2, l}\dfrac{1}{2\varepsilon} + C_{2, l}\dfrac{1}{2\Tilde{\varepsilon}}\right) - 1\right
    )\lVert \nabla\mu_{j}^{k+1} \rVert^{2}_{L^{2}(\Omega_{j})}+\\
    &\quad + \dfrac{\hat{\varepsilon}}{2}\lVert \Delta\eta_{j}^{k+1} \rVert^{2}_{L^{2}(\Omega_{j})} + \dfrac{1}{2\hat{\varepsilon}}\lVert \mu_{j}^{k+1} \rVert^{2}_{L^{2}(\Omega_{j})}.
\end{align*}
Recall that the inverse estimate for finite-element space (see 4.5 Inverse Estimates in “The Mathematical Theory of Finite Element Methods” by Susanne C. Brenner and L. Ridgway Scott) is given by
\begin{equation*}
    \lVert v \rVert_{H^{1}(\Omega)} = \left(\sum_{T\in \mathcal{T}_{h}}\lVert v \rVert^{2}_{H^{1}(T)}\right)^{1/2} \leq Ch^{-1}\left(\sum_{T\in \mathcal{T}_{h}}\lVert v \rVert^{2}_{L^{2}(T)}\right)^{1/2} = Ch^{-1}\lVert v \rVert_{L^{2}(\Omega)}
\end{equation*}
for all $v\in V_{h} = \{v: v \text{ is measurable and } \left.v\right|_{T}\in \mathcal{P}_{T}, \forall T\in \mathcal{T}_{h}\}$, where $\mathcal{P}_{T}\subseteq H^{1}(T)\cap L^{2}(T)$ and $h = \max_{T\in\mathcal{T}_{h}}\text{diam}(T)$, and $C$ depends on the polynomial degree and shape regularity of the mesh. This implies that
\begin{equation*}
    \lVert \nabla v \rVert^{2}_{L^{2}(\Omega)} \leq C^{2}h^{-2}\lVert v \rVert^{2}_{L^{2}(\Omega)}.
\end{equation*}
Since $\mu_{j}^{k+1} \in V_{j,r} \subset V_j$ is a finite-dimensional space of piecewise polynomials on a shape-regular mesh $\mathcal{T}_{h,j}$, the standard inverse inequality (Brenner \& Scott, Thm. 4.5.11) implies
\begin{equation*}
    \|\nabla \mu_{j}^{k+1}\|_{L^2(\Omega_j)} \leq C_{j} h_j^{-1} \|\mu_{j}^{k+1}\|_{L^2(\Omega_j)},
\end{equation*}
where $C$ depends on the polynomial degree and the mesh regularity. This allows us to control the $H^1$-norm of $\mu_{j}^{k+1}$ in terms of its $L^2$-norm. Thus, the inequality becomes 
\begin{align*}
    \dfrac{d}{dt}\left\lVert \mu_{j}^{k+1}\right\rVert^{2}_{L^{2}(\Omega_{j})} &\leq \varepsilon\sum_{l\in \mathcal{N}_{j}}C_{1, l}C_{2, l}\lVert e_{l, r}^{k} \rVert^{2}_{H^{1}(\Omega_{l})} + \Tilde{\varepsilon}\sum_{l\in \mathcal{N}_{j}}C_{2, l}\inf_{v_{r}\in V_{j, r}}\lVert \mathcal{T}_{jl}u_{l}^{k} -  v_{r} \rVert_{H^{1/2}(\Gamma_{jl})}+\\
    &\quad + \left(\dfrac{1}{\hat{\varepsilon}} + C_{j}^{2}h^{-2}_{j}\left(\sum_{l\in \mathcal{N}_{j}}\left(C_{1, l}C_{2, l}\dfrac{1}{\varepsilon} + C_{2, l}\dfrac{1}{\Tilde{\varepsilon}}\right) - 2\right
    )\right)\lVert \mu_{j}^{k+1} \rVert^{2}_{L^{2}(\Omega_{j})}+\\
    &\quad + \hat{\varepsilon}\lVert \Delta\eta_{j}^{k+1} \rVert^{2}_{L^{2}(\Omega_{j})}.
\end{align*}
We can rewrite this as
\begin{equation*}
    \dot{\varphi}_{j}^{k+1}(t) \leq \beta_{j}(t)\varphi(t) + \alpha_{j}^{k}(t),
\end{equation*}
where
\begin{align*}
    \varphi_{j}^{k+1}(t) &= \left\lVert \mu_{j}^{k+1}(t)\right\rVert^{2}_{L^{2}(\Omega_{j})},\\
    \beta_{j}(t) &= \dfrac{1}{\hat{\varepsilon}} + C_{j}^{2}h^{-2}_{j}\left(\sum_{l\in \mathcal{N}_{j}}\left(C_{1, l}C_{2, l}\dfrac{1}{\varepsilon} + C_{2, l}\dfrac{1}{\Tilde{\varepsilon}}\right) - 2\right
    ),\\
    \alpha_{j}^{k}(t) &= \varepsilon\sum_{l\in \mathcal{N}_{j}}C_{1, l}C_{2, l}\lVert e_{l, r}^{k} \rVert^{2}_{H^{1}(\Omega_{l})} + \Tilde{\varepsilon}\sum_{l\in \mathcal{N}_{j}}C_{2, l}\inf_{v_{r}\in V_{j, r}}\lVert \mathcal{T}_{jl}u_{l}^{k} -  v_{r} \rVert^{2}_{H^{1/2}(\Gamma_{jl})} + \hat{\varepsilon}\lVert \Delta\eta_{j}^{k+1} \rVert^{2}_{L^{2}(\Omega_{j})}.
\end{align*}
Then, by Grönwall's lemma, we have
\begin{equation*}
    \varphi_{j}^{k+1}(t) \leq \varphi_{j}^{k+1}(0)e^{\int_{0}^{t}\beta_{j}\, d\tau} + \int_{0}^{t}\alpha_{j}^{k}(s)e^{\int_{s}^{t}\beta_{j}\, d\tau}\,ds.
\end{equation*}
Since $\beta_{j}(t)$ is constant and $\mu_{j}^{k+1}(\cdot, 0) = 0$ for all $k$, we have
\begin{equation*}
    \varphi_{j}^{k+1}(t) \leq \int_{0}^{t}\alpha_{j}^{k}(s)e^{\beta_{j}(t-s)}\,ds \leq \sup_{s\in[0, t]}\alpha_{j}^{k+1}(s)\cdot\int_{0}^{t}e^{\beta_{j}(t-s)}\,ds = \dfrac{e^{\beta_{j} t} - 1}{\beta_{j}}\cdot\sup_{s\in[0, t]}\alpha_{j}^{k}(s) = M_{j}(t)\sup_{s\in[0, t]}\alpha_{j}^{k}(s).
\end{equation*}
Let us take $\varepsilon = \Tilde{\varepsilon} = \hat{\varepsilon} = 1$, then
\begin{align*}
    \varphi_{j}^{k+1}(t) &= \left\lVert \mu_{j}^{k+1}(t)\right\rVert^{2}_{L^{2}(\Omega_{j})},\\
    \beta_{j} &= 1 + C_{j}^{2}h^{-2}_{j}\left(\sum_{l\in \mathcal{N}_{j}}C_{2, l}\left(C_{1, l} + 1\right) - 2\right
    ),\\
    \alpha_{j}^{k}(t) &= \sum_{l\in \mathcal{N}_{j}}C_{1, l}C_{2, l}\lVert e_{l, r}^{k} \rVert^{2}_{H^{1}(\Omega_{l})} + \sum_{l\in \mathcal{N}_{j}}C_{2, l}\inf_{v_{r}\in V_{j, r}}\lVert \mathcal{T}_{jl}u_{l}^{k} -  v_{r} \rVert^{2}_{H^{1/2}(\Gamma_{jl})} + \lVert \Delta\eta_{j}^{k+1} \rVert^{2}_{L^{2}(\Omega_{j})}.
\end{align*}
Then

\begin{align*}
    \left\lVert \mu_{j}^{k+1}(t)\right\rVert^{2}_{L^{2}(\Omega_{j})} &\leq M_{j}(t)\left(\sup_{s\in[0, t]}\sum_{l\in \mathcal{N}_{j}}C_{1, l}C_{2, l}\lVert e_{l, r}^{k} \rVert^{2}_{H^{1}(\Omega_{l})} + \right.\\
    &\quad + \left.\sup_{s\in[0, t]}\sum_{l\in \mathcal{N}_{j}}C_{2, l}\inf_{v_{r}\in V_{j, r}}\lVert \mathcal{T}_{jl}u_{l}^{k} -  v_{r} \rVert^{2}_{H^{1/2}(\Gamma_{jl})} + \sup_{s\in[0, t]}\lVert \Delta\eta_{j}^{k+1} \rVert^{2}_{L^{2}(\Omega_{j})}\right).
\end{align*}

\begin{equation*}
    \sup_{s\in[0, t]}\alpha(s) = \sup_{s\in[0, t]}\sum_{l\in \mathcal{N}_{j}}C_{1, l}C_{2, l}\lVert e_{l, r}^{k} \rVert^{2}_{H^{1}(\Omega_{l})} + \sup_{s\in[0, t]}\sum_{l\in \mathcal{N}_{j}}C_{2, l}\inf_{v_{r}\in V_{j, r}}\lVert \mathcal{T}_{jl}u_{l}^{k} -  v_{r} \rVert^{2}_{H^{1/2}(\Gamma_{jl})} + \sup_{s\in[0, t]}\lVert \Delta\eta_{j}^{k+1} \rVert^{2}_{L^{2}(\Omega_{j})}
\end{equation*}


\begin{equation*}
    \lVert e_{l, r}^{k} \rVert^{2}_{H^{1}(\Omega_{l})} = \lVert \eta_{l}^{k} \rVert^{2}_{H^{1}(\Omega_{l})} + \lVert \mu_{l}^{k} \rVert^{2}_{H^{1}(\Omega_{l})} \leq \lVert \eta_{l}^{k} \rVert^{2}_{H^{1}(\Omega_{l})} + C^{2}_{l}h^{-2}_{l}M_{l}(t)\sup_{s\in[0, t]}\alpha_{l}^{k}(s)
\end{equation*}

\end{document}
